% Candidate D.63; notes only, independent audit before integration.
\subsection{The separated global obstruction and its exact local convergence condition}

\begin{statement}{D.63}{Proved}{A noncompact separated obstruction for every Alexander twist}
Let \(s=(a,b)\), \(A_s=\Lambda C\), \(\mathfrak m\), and
\(W_r=A_s/\mathfrak m^rA_s\) be as in D.62. Arbitrary
infinite ramification is allowed. Give continuous cohomology
the uniform-cochain quotient topology of that statement and put
\[
 \mathcal H_s=H^2_{\mathrm{cts}}(G_{\mathbb Q},A_s)/
                                      \overline{\{0\}}.
                                                        \tag{D.63.1}
\]
For each rational prime \(q\ne5\), choose tame inertia and
an arithmetic Frobenius lift, and define the actual operators
\[
 \begin{split}
 \mu_q&=t_1^{a_q}t_2^{b_q},\qquad
 D_q=q^{-1}\rho_s(\sigma_q)-1,\\
 J_q&=(\mu_q-1)A_s+D_qA_s,\qquad L_q=A_s/J_q,\\
 F^rL_q&=(\mathfrak m^rA_s+J_q)/J_q.
 \end{split}                                             \tag{D.63.2}
\]
The sum defining \(J_q\) is a closed additive subgroup;
\(D_q\) is the actual arithmetic operator of D.35 and D.39,
not a weight approximation. Define
\[
 \mathcal L_s=\left\{(\lambda_q)_{q\ne5}\in\prod_{q\ne5}L_q:
  \text{for every }r,\quad
  \lambda_q\in F^rL_q\text{ for all but finitely many }q\right\}.
                                                        \tag{D.63.3}
\]
Its topology is the inverse-limit topology of the discrete
direct sums \(\bigoplus_{q\ne5}L_q/F^rL_q\).
There is a topologically exact sequence
\[
 0\longrightarrow\mathcal K_s\longrightarrow\mathcal H_s
       \xrightarrow{\mathrm{loc}_{\ne5}}\mathcal L_s
       \longrightarrow0,                               \tag{D.63.4}
\]
where \(\mathcal K_s\) is compact and the displayed map is
open. In particular every family satisfying precisely the
convergence condition (D.63.3) occurs as the localizations
away from five of a global continuous second-cohomology
class.

Here is the exact remaining compact term. Put
\[
 \begin{split}
 T_r&=H^0(G_{\mathbb Q},W_r^D)^*,\qquad
 W_r^D=\operatorname{Hom}(W_r,\mu_{5^\infty}),\\
 \mathcal T_s&=\varprojlim_r T_r,
 \qquad L_5=H^2_{\mathrm{cts}}(G_{\mathbb Q_5},A_s),\\
 \mathcal S_s&=\varprojlim_r
   \ker\left(H^2(G_{\mathbb Q},W_r)
                   \longrightarrow\bigoplus_vH^2(G_{\mathbb Q_v},W_r)\right).
 \end{split}                                             \tag{D.63.5}
\]
The star denotes the finite Pontryagin dual. All terms
in the inverse limit defining \(\mathcal S_s\) are finite.
There are exact sequences of compact groups
\[
 \begin{gathered}
 L_5\longrightarrow\mathcal T_s\longrightarrow0,\\
 0\longrightarrow\mathcal S_s\longrightarrow\mathcal K_s
     \longrightarrow\ker(L_5\longrightarrow\mathcal T_s)
     \longrightarrow0.
 \end{gathered}                                         \tag{D.63.6}
\]
These compact terms are retained, not asserted to vanish.

For every \(s\), \(\mathcal H_s\) is a complete,
countably based, separable Hausdorff group that is not
locally compact. It has a continuous quotient
\[
           \bigoplus_{q\equiv1\ (\mathrm{mod}\ 5)}\mathbb F_5.
                                                        \tag{D.63.7}
\]
Thus D.62's nonseparated kernel is accompanied by a large
separated obstruction, including when \(H^1(G_{\mathbb Q},A_s)=0\).
The compact groups in (D.63.6), and the isomorphism type of
the nonseparated kernel, remain separate unevaluated parts
of Task D.5 for a general twist.
\end{statement}

\begin{proof}
Write \(B_r=H^2(G_{\mathbb Q},W_r)\). The odd-primary
cohomological dimension of \(G_{\mathbb Q}\) is two;
equivalently apply \cite[I, Theorem 4.10(c)]{DE_MilneADT}
with all primes allowed, since the positive-degree
cohomology of a real decomposition group with five-primary
coefficients is zero. The coefficient long exact sequence
therefore makes \(B_{r+1}\to B_r\) surjective. Each
\(B_r\) is countable, by the finite-quotient argument in
D.62. That statement gives an algebraic identification
\[
                    \mathcal H_s\simeq\varprojlim_r B_r.
                                                        \tag{D.63.8}
\]
This is also topological. If a continuous two-cocycle has
zero class modulo \(\mathfrak m^r\), lift a finite-valued
one-cochain witnessing that fact and subtract its
coboundary. The resulting representative takes values
in \(\mathfrak m^rA_s\). Thus the kernel of reduction to
\(B_r\) is exactly the image of the uniform neighborhood
of zero at level \(r\). These kernels define both
topologies in (D.63.8). In particular \(\mathcal H_s\)
is a complete separable metrizable group.

\paragraph{The local quotients at every finite level.}
At \(q\ne5\), the original tripod action is unramified.
Its twisted action on inertia factors through five-primary
tame inertia, whose generator acts by \(\mu_q\).
Wild inertia and the other tame factors have no positive
five-primary cohomology here. The usual procyclic
resolution and the tame relation give
\[
 \begin{split}
 H^1(I_q,W_r)&=W_r/(\mu_q-1)W_r,\\
 H^2(G_{\mathbb Q_q},W_r)
   &=W_r/\bigl((\mu_q-1)W_r+D_qW_r\bigr)\
    =A_s/(J_q+\mathfrak m^rA_s).
 \end{split}                                             \tag{D.63.9}
\]
The Frobenius action on the first line is
\(q^{-1}\rho_s(\sigma_q)\), because the inertia generator
has cyclotomic weight one. Hochschild--Serre has no
other term in degree two: both the relevant inertia and
the residue-field group have five-cohomological dimension
one. This proves (D.63.9) even when \(\mu_q-1\) is a
zero divisor on the finite module, including \(\mu_q=1\).
Changing the Frobenius lift changes \(D_q\) only by an
operator whose image is contained in \((\mu_q-1)A_s\),
so the quotient is independent of that choice.

The two images defining \(J_q\) are compact, so their
sum is compact and closed. Finiteness of local first
cohomology eliminates its first derived inverse limit
\cite[I, Corollary 2.3]{DE_MilneADT}. Taking inverse
limits in (D.63.9) gives
\[
                   H^2_{\mathrm{cts}}(G_{\mathbb Q_q},A_s)=L_q.
                                                        \tag{D.63.10}
\]
In particular each finite local transition is onto,
also directly by its displayed quotient description.
Taking the inverse limit of the direct sums of (D.63.9)
gives exactly (D.63.3): its elements are compatible local
classes with finite support at every finite precision.
There is no common finite ramification set for \(A_s\)
in this argument.

\paragraph{Global reciprocity is absorbed at five.}
Each finite module \(W_r\) has finite ramification,
although the compact module \(A_s\) need not.
Apply finite-module global duality with the set of all
places. The last part of the Poitou--Tate sequence is
\[
 0\longrightarrow \operatorname{Sha}^2(\mathbb Q,W_r)
 \longrightarrow B_r
 \longrightarrow\bigoplus_v H^2(G_{\mathbb Q_v},W_r)
 \longrightarrow T_r\longrightarrow0.
                                                        \tag{D.63.11}
\]
Here \(\operatorname{Sha}^2\) denotes the kernel displayed
in (D.63.5), not an elliptic-curve Tate--Shafarevich group.
It is finite by \cite[I, Theorem 4.10(a)]{DE_MilneADT};
the exactness is part (b). The local summands are finite.
The last map is the sum of the local Tate pairings with
global invariant elements of \(W_r^D\).

The restriction
\(H^0(G_{\mathbb Q},W_r^D)\hookrightarrow
  H^0(G_{\mathbb Q_5},W_r^D)\)
is injective. Finite local duality therefore makes the
single component
\[
                  H^2(G_{\mathbb Q_5},W_r)\longrightarrow T_r
                                                        \tag{D.63.12}
\]
surjective. For any finite-support choice of local
classes away from five, add a class at five cancelling
its image in \(T_r\). Exactness of (D.63.11) then
realizes the resulting family globally. This is the
argument in \cite[I, Corollary 4.16 and its proof,
printed p.62]{DE_MilneADT}; it prescribes zero at all
other places as well, rather than merely prescribing
the originally chosen finite subset.

Let
\[
 K_r=\ker\left(B_r\longrightarrow
                      \bigoplus_{q\ne5}H^2(G_{\mathbb Q_q},W_r)\right).
\]
The real summand is zero. Equations (D.63.11)--(D.63.12)
give exact sequences
\[
 \begin{gathered}
 0\longrightarrow K_r\longrightarrow B_r
   \longrightarrow\bigoplus_{q\ne5}H^2(G_{\mathbb Q_q},W_r)
   \longrightarrow0,\\
 0\longrightarrow\operatorname{Sha}^2(\mathbb Q,W_r)
   \longrightarrow K_r
   \longrightarrow
       \ker\{H^2(G_{\mathbb Q_5},W_r)\longrightarrow T_r\}
   \longrightarrow0.
 \end{gathered}                                          \tag{D.63.13}
\]
In particular \(K_r\) is finite. All inverse systems
of the finite kernels in these rows are Mittag--Leffler.
Taking inverse limits proves (D.63.4) with
\(\mathcal K_s=\varprojlim_rK_r\), and proves (D.63.6).
No surjectivity of the individual transition maps on
the finite kernels is needed.

We also verify the asserted quotient topology. Fix
\(r\), and choose \(t\ge r\) so that the images
of \(K_j\to K_r\) have stabilized for \(j\ge t\).
The stable image equals the image of
\(\varprojlim K_j\to K_r\), by the finite inverse-limit
property. If \(\lambda\in\mathcal L_s\) is zero at
level \(t\), choose a global lift \(h\in\mathcal H_s\).
Then \(h_t\in K_t\), so \(h_r\) is the reduction of
some \(k\in\mathcal K_s\). The lift \(h-k\) is zero
at level \(r\). Thus the image of the \(r\)-th basic
open subgroup of \(\mathcal H_s\) contains the
\(t\)-th basic open subgroup of \(\mathcal L_s\).
This proves openness and topological exactness.

\paragraph{Noncompactness at every scale.}
At level one the action is always
\(W_1=\mathbb F_5(2)\), independent of \(s\).
Formula (D.63.9) gives
\[
 H^2(G_{\mathbb Q_q},W_1)=
 \begin{cases}\mathbb F_5,&q\equiv1\pmod5,\\
 0,&q\not\equiv1\pmod5
 \end{cases}\qquad(q\ne5).
                                                        \tag{D.63.14}
\]
The local transitions and therefore the transitions
between their direct sums are onto. Projection from
their inverse limit is onto, giving (D.63.7).

In fact no basic open subgroup of \(\mathcal L_s\)
is compact. Fix \(r\). Choose infinitely many primes
splitting the actual finite action on \(W_{r+1}\)
and \(\mu_{5^{r+1}}\); they exist by Chebotarev
\cite[Chapter V, Theorem 3.23]{PBC_MilneCFT}.
At each of these primes both terms in the local
differential vanish modulo \(\mathfrak m^{r+1}\),
so the local groups at levels \(r+1\) and \(r\)
are respectively \(W_{r+1}\) and \(W_r\).
Their nonzero kernel is
\(\mathfrak m^rA_s/\mathfrak m^{r+1}A_s\).
An arbitrary finite-support collection of these
kernel classes lifts to a compatible element of
\(\mathcal L_s\), using the surjectivity of its
finite-level transitions. The basic subgroup zero
at level \(r\) consequently has an infinite discrete
quotient at level \(r+1\). It cannot be compact.
Every neighborhood of zero contains one of these
closed open subgroups, so \(\mathcal L_s\) is not
locally compact. The open quotient map in (D.63.4)
then shows that \(\mathcal H_s\) is not locally
compact either. This completes the proof while
retaining exactly the compact ambiguity in (D.63.6).
\end{proof}
